
%package list
\documentclass{article}
\usepackage{amsmath}
\usepackage[top=3cm, bottom=3cm, outer=3cm, inner=3cm]{geometry}
\usepackage{multicol}
\usepackage{graphicx}
\usepackage{url}
%\usepackage{cite}
\usepackage{hyperref}
\hypersetup{
    colorlinks=true,
    allcolors=black,
    linktoc=all,
    pdfborder={0 0 0}
}
\usepackage{array}
%\usepackage{multicol}
\newcolumntype{x}[1]{>{\centering\arraybackslash\hspace{0pt}}p{#1}}
\usepackage[round,authoryear]{natbib}
\setcounter{tocdepth}{2}
\usepackage{pdfpages}
\usepackage{multirow}
\usepackage[normalem]{ulem}
\useunder{\uline}{\ul}{}
\usepackage{svg}
\usepackage{xcolor}
\usepackage{listings}
\lstdefinestyle{ascii-tree}{
    literate={├}{|}1 {─}{--}1 {└}{+}1 
  }
\lstset{basicstyle=\ttfamily,
  showstringspaces=false,
  commentstyle=\color{red},
  keywordstyle=\color{blue}
}
%\usepackage{booktabs}
\usepackage{caption}
\usepackage{subcaption}
\usepackage{float}
\usepackage{array}

\newcolumntype{M}[1]{>{\centering\arraybackslash}m{#1}}
\newcolumntype{N}{@{}m{0pt}@{}}


%%%%%%%%%%%%%%%%%%%%%%%%%%%%%%%%%%%%%%%%%%%%%%%%%%%%%%%%%%%%%%%%%%%%%%%%%%%%
%%%%%%%%%%%%%%%%%%%%%%%%%%%%%%%%%%%%%%%%%%%%%%%%%%%%%%%%%%%%%%%%%%%%%%%%%%%%
\newcommand{\itemEmail}{ccapatinta@unsa.edu.pe}
\newcommand{\itemStudent}{Capatinta Almanza Camilo B. A.}
\newcommand{\itemCourse}{Análisis y Diseño de Algoritmos}
\newcommand{\itemCourseCode}{1702231}
\newcommand{\itemSemester}{IV}
\newcommand{\itemUniversity}{Universidad Nacional de San Agustín de Arequipa}
\newcommand{\itemFaculty}{Facultad de Ingeniería de Producción y Servicios}
\newcommand{\itemDepartment}{Departamento Académico de Ingeniería de Sistemas e Informática}
\newcommand{\itemSchool}{Escuela Profesional de Ingeniería de Sistemas}
\newcommand{\itemAcademic}{2026 - B}
\newcommand{\itemInput}{Del 17 septiembre 2026}
\newcommand{\itemOutput}{Al 20 septiembre 2026}
\newcommand{\itemPracticeNumber}{02}
\newcommand{\itemTheme}{Teorema maestro y relaciones de recurrencia}
%%%%%%%%%%%%%%%%%%%%%%%%%%%%%%%%%%%%%%%%%%%%%%%%%%%%%%%%%%%%%%%%%%%%%%%%%%%%
%%%%%%%%%%%%%%%%%%%%%%%%%%%%%%%%%%%%%%%%%%%%%%%%%%%%%%%%%%%%%%%%%%%%%%%%%%%%

\usepackage[english,spanish]{babel}
\usepackage[utf8]{inputenc}
\AtBeginDocument{\selectlanguage{spanish}}
\renewcommand{\figurename}{Figura}
\renewcommand{\refname}{Referencias}
\renewcommand{\tablename}{Tabla} %esto no funciona cuando se usa babel
\AtBeginDocument{%
	\renewcommand\tablename{Tabla}
}

\usepackage{fancyhdr}
\pagestyle{fancy}
\fancyhf{}
\setlength{\headheight}{40.51407pt}
\renewcommand{\headrulewidth}{1pt}
\renewcommand{\footrulewidth}{1pt}
\fancyhead[L]{\raisebox{-0.2\height}{\includegraphics[width=3cm]{img/logo_episunsa.png}}}
\fancyhead[C]{\fontsize{7}{7}\selectfont	\itemUniversity \\ \itemFaculty \\ \itemDepartment \\ \itemSchool \\ \textbf{\itemCourse}}
\fancyhead[R]{\raisebox{-0.2\height}{\includegraphics[width=1.2cm]{img/logo_abet}}}
\fancyfoot[L]{Estudiante Capatinta Camilo}
\fancyfoot[C]{\itemCourse}
\fancyfoot[R]{Página \thepage}

% para el codigo fuente
\usepackage{listings}
\usepackage{color, colortbl}
\definecolor{dkgreen}{rgb}{0,0.6,0}
\definecolor{gray}{rgb}{0.5,0.5,0.5}
\definecolor{mauve}{rgb}{0.58,0,0.82}
\definecolor{codebackground}{rgb}{0.95, 0.95, 0.92}
\definecolor{tablebackground}{rgb}{0.8, 0, 0}

\lstset{frame=tb,
	language=bash,
	aboveskip=3mm,
	belowskip=3mm,
	showstringspaces=false,
	columns=flexible,
	basicstyle={\small\ttfamily},
	numbers=none,
	numberstyle=\tiny\color{gray},
	keywordstyle=\color{blue},
	commentstyle=\color{dkgreen},
	stringstyle=\color{mauve},
	breaklines=true,
	breakatwhitespace=true,
	tabsize=3,
	backgroundcolor= \color{codebackground},
}

\begin{document}

	\vspace*{10px}
	
	\begin{center}	
		\fontsize{17}{17} \textbf{ Actividad \itemPracticeNumber}
	\end{center}
	\centerline{\textbf{\Large Tema: \itemTheme}}
	%\vspace*{0.5cm}	

	\begin{flushright}
		\begin{tabular}{|M{2.5cm}|N|}
			\hline 
			\rowcolor{tablebackground}
			\color{white} \textbf{Nota}  \\
			\hline 
			     \\[30pt]
			\hline 			
		\end{tabular}
	\end{flushright}	

	\begin{table}[h!]
		\begin{tabular}{|x{4.7cm}|x{4.8cm}|x{4.8cm}|}        
			\hline 
			\rowcolor{tablebackground}
			\color{white} \textbf{Estudiante} & \color{white}\textbf{Escuela}  & \color{white}\textbf{Asignatura}   \\
			\hline 
			{\itemStudent \par \itemEmail} & \itemSchool & {\itemCourse \par Semestre: \itemSemester \par Código: \itemCourseCode}     \\
			\hline 			
		\end{tabular}
	\end{table}		
	
	\begin{table}[H]
		\begin{tabular}{|x{4.7cm}|x{4.8cm}|x{4.8cm}|}
			\hline 
			\rowcolor{tablebackground}
			\color{white}\textbf{Actividad} & \color{white}\textbf{Tema}  & \color{white}\textbf{Duración}   \\
			\hline 
			\itemPracticeNumber & \itemTheme & 04 horas   \\
			\hline 
		\end{tabular}
	\end{table}
	
	\begin{table}[H]
		\begin{tabular}{|x{4.7cm}|x{4.8cm}|x{4.8cm}|}
			\hline 
			\rowcolor{tablebackground}
			\color{white}\textbf{Semestre académico} & \color{white}\textbf{Fecha de inicio}  & \color{white}\textbf{Fecha de entrega}   \\
			\hline 
			\itemAcademic & \itemInput &  \itemOutput  \\
			\hline 
		\end{tabular}
	\end{table}
\clearpage
    \color{black}
\renewcommand{\contentsname}{\textbf{Índice}}
\tableofcontents
\clearpage

    \section{Teorema Maestro y Relaciones de Recurrencia}

    Investigar el Teorema Maestro como herramienta para resolver relaciones de
    recurrencia de la forma $T(n) = a\,T(n/b) + f(n)$, con $a \geq 1$ y $b > 1$.

        \begin{itemize}
        \item Indicar los casos de análisis que contempla el teorema, según la
                comparación entre $f(n)$ y $n^{\log_b a}$.
        \item Proponer tres casos de uso, planteando la relación de recurrencia de
                cada uno y resolviéndola con el teorema.
        \end{itemize}
		
	\section{Teorema maestro}
 
        \subsection{Objetivo}
        
        Investigar el Teorema Maestro como herramienta para resolver relaciones de
        recurrencia de la forma
        \[
        T(n) = a\,T\!\left(\frac{n}{b}\right) + f(n),
        \]
        indicar los casos de análisis que contempla según la comparación entre $f(n)$
        y $n^{\log_b a}$, y proponer tres casos de uso, planteando la relación de
        recurrencia de cada uno y resolviéndola con el teorema.
        
        % ------------------------------------------------------------
        \subsection{Relaciones de recurrencia y divide y vencerás}
        
        Muchos algoritmos siguen el paradigma de \emph{divide y vencerás}: dividen el
        problema en subproblemas más pequeños, los resuelven de forma recursiva y
        combinan sus soluciones. El tiempo de ejecución de estos algoritmos se expresa
        de manera natural mediante una \textbf{relación de recurrencia}, es decir, una
        ecuación que define $T(n)$ en términos del costo sobre entradas más pequeñas.
        El Teorema Maestro proporciona una solución asintótica directa para las
        recurrencias de este tipo sin necesidad de desarrollarlas manualmente
        \citep{clrs}.
        
        % ------------------------------------------------------------
        \subsection{Forma general de la recurrencia}
        
        El Teorema Maestro se aplica a recurrencias de la forma
        \[
        T(n) = a\,T\!\left(\frac{n}{b}\right) + f(n),
        \qquad a \geq 1,\; b > 1,
        \]
        donde:
        
        \begin{itemize}
        \item $n$ es el tamaño del problema.
        \item $a$ es el número de subproblemas en que se divide el problema en cada
                llamada recursiva (constante, $a \geq 1$).
        \item $n/b$ es el tamaño de cada subproblema (constante $b > 1$). Se admite
                que sea $\lfloor n/b \rfloor$ o $\lceil n/b \rceil$ sin alterar el
                resultado asintótico.
        \item $f(n)$ es el costo de dividir el problema y combinar las soluciones de
                los subproblemas; debe ser una función asintóticamente positiva.
        \item Se asume un caso base $T(n) = \Theta(1)$ para $n$ suficientemente
                pequeño.
        \end{itemize}
        
        % ------------------------------------------------------------
        \subsection{Intuición: el árbol de recursión}
        
        Al desarrollar la recurrencia se obtiene un árbol de recursión con las
        siguientes propiedades:
        
        \begin{itemize}
        \item En el nivel $i$ hay $a^i$ subproblemas, cada uno de tamaño $n/b^i$,
                por lo que el costo total de ese nivel es $a^i f(n/b^i)$.
        \item La altura del árbol es $\log_b n$, ya que $n/b^i = 1$ cuando
                $i = \log_b n$.
        \item El número de hojas es $a^{\log_b n} = n^{\log_b a}$, y su costo total
                es $\Theta\!\left(n^{\log_b a}\right)$.
        \end{itemize}
        
        Por lo tanto,
        \[
        T(n) = \Theta\!\left(n^{\log_b a}\right)
            + \sum_{i=0}^{\log_b n - 1} a^i\, f\!\left(\frac{n}{b^i}\right).
        \]
        El primer término es el costo de las hojas y la suma es el costo de los nodos
        internos. El Teorema Maestro determina cuál de los dos domina comparando
        $f(n)$ con la \textbf{función crítica} $n^{\log_b a}$.
        
        % ------------------------------------------------------------
        \subsection{Enunciado del Teorema Maestro}
        
        \begin{quote}
        \textbf{Teorema (Maestro).} Sean $a \geq 1$ y $b > 1$ constantes, sea $f(n)$
        una función asintóticamente positiva y sea $T(n)$ definida para los enteros no
        negativos por la recurrencia $T(n) = a\,T(n/b) + f(n)$. Entonces $T(n)$ tiene
        las siguientes cotas asintóticas:
        \end{quote}
        
        \begin{itemize}
        \item \textbf{Caso 1 (domina el costo de las hojas).} Si
                $f(n) = O\!\left(n^{\log_b a - \varepsilon}\right)$ para alguna
                constante $\varepsilon > 0$, entonces
                \[
                T(n) = \Theta\!\left(n^{\log_b a}\right).
                \]
        
        \item \textbf{Caso 2 (el costo se reparte entre todos los niveles).} Si
                $f(n) = \Theta\!\left(n^{\log_b a}\log^{k} n\right)$ para alguna
                constante $k \geq 0$, entonces
                \[
                T(n) = \Theta\!\left(n^{\log_b a}\log^{k+1} n\right).
                \]
                En particular, para $k = 0$: si
                $f(n) = \Theta\!\left(n^{\log_b a}\right)$, entonces
                $T(n) = \Theta\!\left(n^{\log_b a}\log n\right)$.
        
        \item \textbf{Caso 3 (domina el costo de la raíz).} Si
                $f(n) = \Omega\!\left(n^{\log_b a + \varepsilon}\right)$ para alguna
                constante $\varepsilon > 0$, y si además se cumple la
                \emph{condición de regularidad}
                \[
                a\,f\!\left(\frac{n}{b}\right) \leq c\,f(n)
                \quad \text{para alguna constante } c < 1
                \text{ y } n \text{ suficientemente grande},
                \]
                entonces
                \[
                T(n) = \Theta\!\left(f(n)\right).
                \]
        \end{itemize}
        
        \subsubsection{Resumen de los casos}
        
        \begin{center}
        \renewcommand{\arraystretch}{1.5}
        \begin{tabular}{|c|c|c|}
        \hline
        \textbf{Caso} & \textbf{Comparación de $f(n)$ con $n^{\log_b a}$} & \textbf{Solución} \\
        \hline
        1 & $f(n)$ crece polinomialmente más lento & $\Theta\!\left(n^{\log_b a}\right)$ \\
        \hline
        2 & $f(n)$ crece al mismo ritmo (salvo factores $\log^k n$) & $\Theta\!\left(n^{\log_b a}\log^{k+1} n\right)$ \\
        \hline
        3 & $f(n)$ crece polinomialmente más rápido (y es regular) & $\Theta\!\left(f(n)\right)$ \\
        \hline
        \end{tabular}
        \end{center}
        
        % ------------------------------------------------------------
        \subsection{Procedimiento para aplicar el teorema}
        
        \begin{itemize}
        \item \textbf{Paso 1:} identificar los valores de $a$, $b$ y $f(n)$ en la
                recurrencia.
        \item \textbf{Paso 2:} calcular la función crítica $n^{\log_b a}$.
        \item \textbf{Paso 3:} comparar $f(n)$ con $n^{\log_b a}$ y determinar en
                cuál de los tres casos se encuentra.
        \item \textbf{Paso 4:} si es el caso 3, verificar la condición de
                regularidad.
        \item \textbf{Paso 5:} escribir la solución asintótica correspondiente.
        \end{itemize}
        
        % ------------------------------------------------------------
        \subsection{Limitaciones del teorema}
        
        El Teorema Maestro no cubre todas las recurrencias de divide y vencerás. No es
        aplicable, entre otros, en los siguientes casos:
        
        \begin{itemize}
        \item Cuando $f(n)$ es menor que $n^{\log_b a}$ pero no de forma polinomial
                (existe una brecha entre los casos 1 y 2). Por ejemplo,
                $T(n) = 2T(n/2) + n/\log n$: aquí $n^{\log_2 2} = n$, y $f(n)$ no es
                $O(n^{1-\varepsilon})$ para ningún $\varepsilon > 0$, ni tampoco
                $\Theta(n\log^k n)$ con $k \geq 0$.
        \item Cuando $f(n)$ es mayor que $n^{\log_b a}$ pero no de forma polinomial,
                o cuando no cumple la condición de regularidad (brecha entre los
                casos 2 y 3).
        \item Cuando $a$ no es constante o es menor que $1$, o cuando los
                subproblemas no tienen el mismo tamaño (por ejemplo,
                $T(n) = T(n/3) + T(2n/3) + n$).
        \item Cuando la recurrencia no es de la forma $T(n/b)$, por ejemplo las
                recurrencias por sustracción como $T(n) = T(n-1) + n$.
        \end{itemize}
        
        En esos casos se recurre a otros métodos: sustitución, árbol de recursión o
        el método de Akra--Bazzi \citep{akrabazzi1998}, que generaliza el Teorema
        Maestro. La versión del teorema con divisiones de la forma $a\,T(n/b)+f(n)$ se
        remonta al trabajo de Bentley, Haken y Saxe \citep{bentley1980}.
        
        % ------------------------------------------------------------
        \subsection{Tres casos de uso}
        
        \subsubsection{Caso de uso 1: Búsqueda binaria}
        
        La búsqueda binaria compara el elemento buscado con el elemento central de un
        arreglo ordenado y continúa en una sola de las dos mitades. La comparación y el
        cálculo del punto medio cuestan tiempo constante.
        
        \textbf{Relación de recurrencia:}
        \[
        T(n) = T\!\left(\frac{n}{2}\right) + \Theta(1).
        \]
        
        \textbf{Resolución con el Teorema Maestro:}
        
        \begin{itemize}
        \item Parámetros: $a = 1$, $b = 2$, $f(n) = \Theta(1)$.
        \item Función crítica: $n^{\log_2 1} = n^{0} = 1$.
        \item Comparación: $f(n) = \Theta(1) = \Theta\!\left(n^{0}\log^{0} n\right)$,
                por lo que se cumple el \textbf{caso 2} con $k = 0$.
        \item Solución:
                \[
                T(n) = \Theta\!\left(n^{0}\log n\right) = \Theta(\log n).
                \]
        \end{itemize}
        
        \subsubsection{Caso de uso 2: Merge Sort}
        
        Merge Sort divide el arreglo en dos mitades, ordena cada una de forma
        recursiva y las mezcla en tiempo lineal.
        
        \textbf{Relación de recurrencia:}
        \[
        T(n) = 2\,T\!\left(\frac{n}{2}\right) + \Theta(n).
        \]
        
        \textbf{Resolución con el Teorema Maestro:}
        
        \begin{itemize}
        \item Parámetros: $a = 2$, $b = 2$, $f(n) = \Theta(n)$.
        \item Función crítica: $n^{\log_2 2} = n$.
        \item Comparación: $f(n) = \Theta(n) = \Theta\!\left(n^{\log_2 2}\right)$,
                por lo que se cumple el \textbf{caso 2} con $k = 0$.
        \item Solución:
                \[
                T(n) = \Theta(n\log n).
                \]
        \end{itemize}
        
        \subsubsection{Caso de uso 3: Multiplicación de Karatsuba}
        
        El algoritmo de Karatsuba multiplica dos números de $n$ dígitos reduciendo las
        cuatro multiplicaciones de tamaño $n/2$ del método clásico a solamente tres,
        a costa de sumas y restas adicionales de costo lineal.
        
        \textbf{Relación de recurrencia:}
        \[
        T(n) = 3\,T\!\left(\frac{n}{2}\right) + \Theta(n).
        \]
        
        \textbf{Resolución con el Teorema Maestro:}
        
        \begin{itemize}
        \item Parámetros: $a = 3$, $b = 2$, $f(n) = \Theta(n)$.
        \item Función crítica: $n^{\log_2 3} \approx n^{1.585}$.
        \item Comparación: $f(n) = n = O\!\left(n^{\log_2 3 - \varepsilon}\right)$
                con, por ejemplo, $\varepsilon = 0.5$, ya que $n^{1.085}$ crece más
                rápido que $n$. Se cumple el \textbf{caso 1}.
        \item Solución:
                \[
                T(n) = \Theta\!\left(n^{\log_2 3}\right) \approx \Theta\!\left(n^{1.585}\right),
                \]
                que mejora asintóticamente al $\Theta(n^2)$ del método clásico.
        \end{itemize}
        
        % ------------------------------------------------------------
        \subsection{Ejemplo adicional: el caso 3}
        
        Los tres casos de uso anteriores ilustran los casos 1 y 2. Para el caso 3 se
        puede considerar la selección de un elemento por orden (\emph{Quickselect}) en
        su escenario favorable, donde cada partición divide el arreglo en mitades y
        solo se continúa en una de ellas. La partición cuesta tiempo lineal.
        
        \textbf{Relación de recurrencia:}
        \[
        T(n) = T\!\left(\frac{n}{2}\right) + \Theta(n).
        \]
        
        \textbf{Resolución con el Teorema Maestro:}
        
        \begin{itemize}
        \item Parámetros: $a = 1$, $b = 2$, $f(n) = \Theta(n)$.
        \item Función crítica: $n^{\log_2 1} = 1$.
        \item Comparación: $f(n) = n = \Omega\!\left(n^{0+\varepsilon}\right)$ con
                $\varepsilon = 1$, por lo que corresponde al \textbf{caso 3}.
        \item Regularidad: $a\,f(n/b) = \dfrac{n}{2} \leq \dfrac{1}{2}\,n$, y se
                cumple con $c = 1/2 < 1$.
        \item Solución:
                \[
                T(n) = \Theta(f(n)) = \Theta(n).
                \]
        \end{itemize}
        
        Otro ejemplo que muestra la versión extendida del caso 2 es
        $T(n) = 2T(n/2) + n\log n$: con $a = b = 2$ se tiene
        $f(n) = \Theta(n\log^{1} n)$, es decir $k = 1$, y por tanto
        $T(n) = \Theta(n\log^{2} n)$.
        
        % ------------------------------------------------------------
        \subsection{Resumen de los casos de uso}
        
        \begin{center}
        \small
        \renewcommand{\arraystretch}{1.5}
        \begin{tabular}{|l|c|c|c|c|}
        \hline
        \textbf{Algoritmo} & \textbf{Recurrencia} & $n^{\log_b a}$ & \textbf{Caso} & \textbf{Solución} \\
        \hline
        Búsqueda binaria & $T(n) = T(n/2) + \Theta(1)$ & $1$ & 2 & $\Theta(\log n)$ \\
        \hline
        Merge Sort & $T(n) = 2T(n/2) + \Theta(n)$ & $n$ & 2 & $\Theta(n\log n)$ \\
        \hline
        Karatsuba & $T(n) = 3T(n/2) + \Theta(n)$ & $n^{1.585}$ & 1 & $\Theta(n^{1.585})$ \\
        \hline
        Quickselect (favorable) & $T(n) = T(n/2) + \Theta(n)$ & $1$ & 3 & $\Theta(n)$ \\
        \hline
        \end{tabular}
        \end{center}



	\section{Relaciones de Recurrencia}
 
        \subsection{Objetivo}
        
        Investigar las relaciones de recurrencia como herramienta para describir el
        costo de los algoritmos recursivos, presentar los métodos para resolverlas
        (iteración, sustitución, árbol de recursión, ecuación característica y
        Teorema Maestro) y desarrollar casos de uso planteando y resolviendo la
        recurrencia de cada uno.
        
        % ------------------------------------------------------------
        \subsection{Definición}
        
        Una \textbf{relación de recurrencia} es una ecuación que define cada término de
        una sucesión $a_0, a_1, a_2, \ldots$ en función de uno o más términos
        anteriores. Para que la sucesión quede completamente determinada, la relación
        debe acompañarse de \textbf{condiciones iniciales} (o \textbf{caso base})
        \citep{rosen2019}.
        
        \textbf{Ejemplos clásicos:}
        
        \begin{itemize}
        \item Factorial: $n! = n\cdot(n-1)!$, con $0! = 1$.
        \item Sucesión de Fibonacci: $F_n = F_{n-1} + F_{n-2}$, con $F_0 = 0$ y
                $F_1 = 1$.
        \item Torres de Hanoi: $H_n = 2H_{n-1} + 1$, con $H_0 = 0$.
        \end{itemize}
        
        \textbf{Resolver} una relación de recurrencia significa encontrar una
        \emph{fórmula cerrada} para $a_n$ que no dependa de términos anteriores, o al
        menos una cota asintótica ($O$, $\Omega$, $\Theta$) de su crecimiento.
        
        % ------------------------------------------------------------
        \subsection{Recurrencias en el análisis de algoritmos}
        
        Si $T(n)$ es el tiempo de ejecución de un algoritmo recursivo sobre una entrada
        de tamaño $n$, la recurrencia se plantea a partir de la estructura del
        algoritmo \citep{clrs}:
        
        \begin{itemize}
        \item \textbf{Caso base:} el costo cuando la entrada es lo bastante pequeña
                para resolverse sin recursión (normalmente $T(1) = \Theta(1)$).
        \item \textbf{Llamadas recursivas:} cuántas se realizan y sobre qué tamaño
                de entrada, lo que produce términos como $T(n-1)$ o $a\,T(n/b)$.
        \item \textbf{Costo no recursivo:} el trabajo que se hace fuera de las
                llamadas (dividir, combinar, comparar), expresado como $f(n)$.
        \end{itemize}
        
        Con estos elementos, la recurrencia tiene la forma
        $T(n) = (\text{costo de las llamadas recursivas}) + f(n)$.
        
        % ------------------------------------------------------------
        \subsection{Clasificación de las recurrencias}
        
        \begin{itemize}
        \item \textbf{Según el orden:} el orden es la distancia entre el término
                más reciente y el más antiguo que aparece. $a_n = 2a_{n-1}+1$ es de
                orden 1; $F_n = F_{n-1}+F_{n-2}$ es de orden 2.
        \item \textbf{Lineales y no lineales:} son lineales si los términos previos
                aparecen solo multiplicados por coeficientes y sumados. Ejemplo no
                lineal: $a_n = a_{n-1}^2 + 1$.
        \item \textbf{Homogéneas y no homogéneas:} una recurrencia lineal es
                homogénea si no tiene términos independientes de la sucesión (por
                ejemplo $F_n = F_{n-1} + F_{n-2}$), y no homogénea en caso contrario
                (por ejemplo $H_n = 2H_{n-1} + 1$).
        \item \textbf{Coeficientes constantes o variables:} los coeficientes que
                multiplican a los términos previos son constantes, como en Fibonacci,
                o dependen de $n$.
        \item \textbf{Por sustracción o por división:} en las primeras el tamaño
                disminuye restando una constante, $T(n) = aT(n-c) + f(n)$; en las
                segundas se divide entre una constante, $T(n) = aT(n/b) + f(n)$
                (divide y vencerás).
        \end{itemize}
        
        % ------------------------------------------------------------
        \subsection{Métodos de resolución}
        
        \subsubsection{Método de iteración (sustitución hacia atrás)}
        
        Consiste en \emph{desenrollar} la recurrencia sustituyendo repetidamente el
        término anterior hasta detectar un patrón en la $k$-ésima sustitución y luego
        elegir $k$ para llegar al caso base.
        
        \begin{itemize}
        \item \textbf{Ventaja:} no requiere adivinar la solución.
        \item \textbf{Limitación:} el álgebra puede volverse complicada y el patrón
                general debe demostrarse (normalmente por inducción).
        \end{itemize}
        
        \subsubsection{Método de sustitución (adivinar y demostrar)}
        
        \begin{itemize}
        \item \textbf{Paso 1:} adivinar la forma de la solución (por ejemplo,
                $T(n) \leq c\,n\log n$).
        \item \textbf{Paso 2:} demostrar por inducción que la cota se cumple,
                suponiendo que vale para tamaños menores.
        \item \textbf{Paso 3:} elegir las constantes de forma que también se cumpla
                el caso base.
        \end{itemize}
        
        \subsubsection{Método del árbol de recursión}
        
        Se dibuja el árbol de llamadas recursivas, se calcula el costo de cada nivel y
        se suman los costos de todos los niveles. Sirve para generar una buena
        conjetura que luego puede confirmarse con el método de sustitución.
        
        \subsubsection{Ecuación característica (recurrencias lineales homogéneas de coeficientes constantes)}
        
        Dada la recurrencia
        \[
        a_n = c_1 a_{n-1} + c_2 a_{n-2} + \cdots + c_k a_{n-k},
        \]
        su \textbf{ecuación característica} es
        \[
        r^k - c_1 r^{k-1} - c_2 r^{k-2} - \cdots - c_k = 0.
        \]
        
        \begin{itemize}
        \item \textbf{Raíces distintas} $r_1, \ldots, r_k$: la solución general es
                \[
                a_n = \alpha_1 r_1^{\,n} + \alpha_2 r_2^{\,n} + \cdots + \alpha_k r_k^{\,n}.
                \]
        \item \textbf{Raíz repetida} $r$ de multiplicidad $m$: aporta a la solución
                el término
                \[
                \left(\alpha_0 + \alpha_1 n + \cdots + \alpha_{m-1} n^{m-1}\right) r^{\,n}.
                \]
        \item Las constantes $\alpha_i$ se determinan resolviendo el sistema que
                resulta de aplicar las condiciones iniciales.
        \item \textbf{Caso no homogéneo:} la solución es la suma de la solución
                general de la parte homogénea y de una solución particular de la
                recurrencia completa. Las constantes se ajustan al final con las
                condiciones iniciales.
        \end{itemize}
        
        \subsubsection{Teorema Maestro}
        
        Para recurrencias por división de la forma $T(n) = aT(n/b) + f(n)$ con
        $a \geq 1$ y $b > 1$, el Teorema Maestro da directamente la cota asintótica
        comparando $f(n)$ con $n^{\log_b a}$ (véase la sección anterior).
        
        \subsubsection{Otros métodos}
        
        Las \emph{funciones generadoras} permiten resolver recurrencias que las técnicas
        anteriores no cubren, y el método de \emph{Akra--Bazzi} generaliza el Teorema
        Maestro a subproblemas de tamaños distintos.
        
        % ------------------------------------------------------------
        \subsection{Casos de uso}
        
        \subsubsection{Caso de uso 1: Torres de Hanoi (recurrencia por sustracción)}
        
        Para mover $n$ discos de una varilla a otra, se mueven $n-1$ discos a la
        varilla auxiliar, se mueve el disco mayor y se mueven de nuevo los $n-1$
        discos sobre él \citep{levitin2012}.
        
        \textbf{Relación de recurrencia:}
        \[
        H(n) = 2\,H(n-1) + 1, \qquad H(0) = 0.
        \]
        
        \textbf{Resolución por iteración:}
        \begin{align*}
        H(n) &= 2H(n-1) + 1 \\
            &= 2\bigl(2H(n-2) + 1\bigr) + 1 = 2^2 H(n-2) + 2 + 1 \\
            &= 2^3 H(n-3) + 2^2 + 2 + 1 \\
            &\;\;\vdots \\
            &= 2^k H(n-k) + \left(2^{k-1} + \cdots + 2 + 1\right) = 2^k H(n-k) + 2^k - 1.
        \end{align*}
        Con $k = n$ se llega al caso base $H(0) = 0$:
        \[
        H(n) = 2^n \cdot 0 + 2^n - 1 = 2^n - 1 = \Theta(2^n).
        \]
        
        \textbf{Verificación:} $H(1) = 1$, $H(2) = 3$, $H(3) = 7$, que coinciden con
        $2^n - 1$.
        
        \textbf{Resolución alternativa (recurrencia no homogénea):} la parte homogénea
        $H(n) = 2H(n-1)$ tiene como solución $\alpha 2^n$. Una solución particular
        constante $p$ debe cumplir $p = 2p + 1$, es decir $p = -1$. Entonces
        $H(n) = \alpha 2^n - 1$, y con $H(0) = 0$ se obtiene $\alpha = 1$, lo que
        confirma $H(n) = 2^n - 1$.
        
        \subsubsection{Caso de uso 2: Sucesión de Fibonacci (ecuación característica)}
        
        \textbf{Relación de recurrencia:}
        \[
        F(n) = F(n-1) + F(n-2), \qquad F(0) = 0,\; F(1) = 1.
        \]
        
        \textbf{Resolución con la ecuación característica:}
        
        \begin{itemize}
        \item Ecuación característica: $r^2 - r - 1 = 0$.
        \item Raíces:
                \[
                r_1 = \varphi = \frac{1+\sqrt{5}}{2} \approx 1.618,
                \qquad
                r_2 = \psi = \frac{1-\sqrt{5}}{2} \approx -0.618.
                \]
        \item Solución general: $F(n) = \alpha\,\varphi^{\,n} + \beta\,\psi^{\,n}$.
        \item Condiciones iniciales: de $F(0) = 0$ resulta $\alpha + \beta = 0$, y de
                $F(1) = 1$ resulta $\alpha\varphi + \beta\psi = 1$. Resolviendo,
                $\alpha = \dfrac{1}{\sqrt{5}}$ y $\beta = -\dfrac{1}{\sqrt{5}}$.
        \item Fórmula cerrada (fórmula de Binet):
                \[
                F(n) = \frac{\varphi^{\,n} - \psi^{\,n}}{\sqrt{5}}.
                \]
        \end{itemize}
        
        Como $|\psi| < 1$, el segundo término tiende a $0$ y
        $F(n) = \Theta\!\left(\varphi^{\,n}\right)$. Por eso el algoritmo recursivo
        ingenuo, cuyo tiempo satisface la misma recurrencia más un costo constante,
        tiene crecimiento exponencial.
        
        \subsubsection{Caso de uso 3: Merge Sort (recurrencia por división)}
        
        Merge Sort divide el arreglo en dos mitades, las ordena recursivamente y las
        mezcla en tiempo lineal.
        
        \textbf{Relación de recurrencia:}
        \[
        T(n) = 2\,T\!\left(\frac{n}{2}\right) + c\,n, \qquad T(1) = \Theta(1).
        \]
        
        \textbf{Resolución por iteración} (con $n = 2^k$):
        \begin{align*}
        T(n) &= 2T(n/2) + cn \\
            &= 2\bigl(2T(n/4) + cn/2\bigr) + cn = 4T(n/4) + 2cn \\
            &= 8T(n/8) + 3cn \\
            &\;\;\vdots \\
            &= 2^k T(n/2^k) + k\,c\,n.
        \end{align*}
        Con $k = \log_2 n$ se llega a $T(1)$:
        \[
        T(n) = n\,T(1) + c\,n\log_2 n = \Theta(n\log n).
        \]
        
        \textbf{Comprobación con el Teorema Maestro:} $a = b = 2$,
        $n^{\log_2 2} = n = \Theta(f(n))$, es decir, caso 2 con $k = 0$, que también
        da $\Theta(n\log n)$.
        
        % ------------------------------------------------------------
        \subsection{Ejemplos adicionales por método}
        
        \subsubsection{Método de sustitución}
        
        Sea $T(n) = 2T(\lfloor n/2\rfloor) + n$ con $T(1) = 1$. Se conjetura
        $T(n) \leq c\,n\log_2 n$ para $n \geq 2$.
        
        \begin{itemize}
        \item \textbf{Paso inductivo:} suponiendo la cota para $\lfloor n/2\rfloor$,
                \begin{align*}
                T(n) &\leq 2c\left\lfloor\frac{n}{2}\right\rfloor\log_2\left\lfloor\frac{n}{2}\right\rfloor + n
                        \leq c\,n\log_2\frac{n}{2} + n \\
                    &= c\,n\log_2 n - c\,n + n \leq c\,n\log_2 n
                    \quad \text{si } c \geq 1.
                \end{align*}
        \item \textbf{Casos base:} se toman $n = 2$ y $n = 3$, ya que para $n = 1$
                se tiene $\log_2 1 = 0$ y la cota fallaría. Como $T(2) = 4$ y
                $T(3) = 5$, basta $c = 2$.
        \item \textbf{Conclusión:} $T(n) = O(n\log n)$.
        \end{itemize}
        
        \subsubsection{Método del árbol de recursión}
        
        Sea $T(n) = 3T(n/4) + c\,n^2$.
        
        \begin{itemize}
        \item En el nivel $i$ hay $3^i$ nodos, cada uno de costo $c\,(n/4^i)^2$, por
                lo que el costo del nivel es $\left(\tfrac{3}{16}\right)^{i} c\,n^2$.
        \item La altura es $\log_4 n$ y el número de hojas es
                $3^{\log_4 n} = n^{\log_4 3}$.
        \item La suma de los niveles internos es una serie geométrica decreciente:
                \[
                \sum_{i=0}^{\log_4 n - 1}\left(\frac{3}{16}\right)^{i} c\,n^2
                < c\,n^2\sum_{i=0}^{\infty}\left(\frac{3}{16}\right)^{i}
                = \frac{16}{13}\,c\,n^2.
                \]
        \item El costo de las hojas es $\Theta\!\left(n^{\log_4 3}\right) = o(n^2)$,
                y el costo de la raíz ya es $c\,n^2$. Por tanto, $T(n) = \Theta(n^2)$.
        \end{itemize}
        
        Este resultado coincide con el caso 3 del Teorema Maestro
        ($a=3$, $b=4$, $n^{\log_4 3}\approx n^{0.79}$, y la regularidad se cumple
        porque $3(n/4)^2 = \tfrac{3}{16}\,n^2$, es decir, con $c = 3/16 < 1$).
        
        \subsubsection{Ecuación característica con raíces distintas}
        
        Sea $a_n = 5a_{n-1} - 6a_{n-2}$ con $a_0 = 2$ y $a_1 = 5$.
        
        \begin{itemize}
        \item Ecuación característica: $r^2 - 5r + 6 = 0$, con raíces $r = 2$ y
                $r = 3$.
        \item Solución general: $a_n = \alpha\,2^n + \beta\,3^n$.
        \item Condiciones iniciales: $\alpha + \beta = 2$ y $2\alpha + 3\beta = 5$,
                de donde $\alpha = \beta = 1$.
        \item Solución: $a_n = 2^n + 3^n$. Verificación: $a_2 = 5\cdot 5 - 6\cdot 2 = 13 = 2^2 + 3^2$.
        \end{itemize}
        
        \subsubsection{Ecuación característica con raíz repetida}
        
        Sea $a_n = 4a_{n-1} - 4a_{n-2}$ con $a_0 = 1$ y $a_1 = 4$.
        
        \begin{itemize}
        \item Ecuación característica: $r^2 - 4r + 4 = (r-2)^2 = 0$, con raíz doble
                $r = 2$.
        \item Solución general: $a_n = (\alpha + \beta n)\,2^n$.
        \item Condiciones iniciales: $\alpha = 1$ y $(\alpha + \beta)\cdot 2 = 4$, de
                donde $\beta = 1$.
        \item Solución: $a_n = (n+1)\,2^n$. Verificación: $a_2 = 4\cdot 4 - 4\cdot 1 = 12 = 3\cdot 2^2$.
        \end{itemize}
        
        % ------------------------------------------------------------
        \subsection{Resumen}
        
        \begin{center}
        \small
        \renewcommand{\arraystretch}{1.5}
        \begin{tabular}{|l|l|c|}
        \hline
        \textbf{Recurrencia} & \textbf{Método habitual} & \textbf{Solución} \\
        \hline
        $T(n) = T(n-1) + c$ & Iteración & $\Theta(n)$ \\
        \hline
        $T(n) = T(n-1) + n$ & Iteración & $\Theta(n^2)$ \\
        \hline
        $T(n) = 2T(n-1) + 1$ & Iteración / no homogénea & $\Theta(2^n)$ \\
        \hline
        $T(n) = T(n-1) + T(n-2)$ & Ecuación característica & $\Theta(\varphi^{\,n})$ \\
        \hline
        $T(n) = T(n/2) + c$ & Iteración / Teorema Maestro & $\Theta(\log n)$ \\
        \hline
        $T(n) = 2T(n/2) + n$ & Iteración / Teorema Maestro & $\Theta(n\log n)$ \\
        \hline
        $T(n) = 3T(n/4) + n^2$ & Árbol / Teorema Maestro & $\Theta(n^2)$ \\
        \hline
        \end{tabular}
        \end{center}
        
        \textbf{Cuándo usar cada método:}
        
        \begin{itemize}
        \item \textbf{Iteración:} recurrencias simples por sustracción o por división
                cuando se puede detectar el patrón.
        \item \textbf{Árbol de recursión:} para obtener una conjetura o cuando el
                costo por nivel forma una serie reconocible.
        \item \textbf{Sustitución:} para demostrar formalmente una cota conjeturada.
        \item \textbf{Ecuación característica:} recurrencias lineales de coeficientes
                constantes, como Fibonacci.
        \item \textbf{Teorema Maestro:} recurrencias de la forma $aT(n/b) + f(n)$ que
                cumplan sus hipótesis.
        \end{itemize}


    

	\bibliographystyle{plainnat}
    \bibliography{referencias}
%\clearpage
%\bibliographystyle{apalike}
%\bibliographystyle{IEEEtranN}
%\bibliography{bibliography}
			
\end{document}